\documentclass[11pt,a4paper]{article}

\usepackage[margin=2.5cm]{geometry}
\usepackage{amsmath,amssymb,amsthm}
\usepackage{booktabs}
\usepackage{array}
\usepackage{hyperref}
\usepackage[most]{tcolorbox}
\usepackage{multirow}

\theoremstyle{plain}
\newtheorem{theorem}{Theorem}[section]
\newtheorem{proposition}[theorem]{Proposition}
\newtheorem{corollary}[theorem]{Corollary}
\newtheorem{lemma}[theorem]{Lemma}
\newtheorem{conjecture}[theorem]{Conjecture}
\theoremstyle{definition}
\newtheorem{definition}[theorem]{Definition}
\newtheorem{observation}[theorem]{Observation}
\newtheorem{remark}[theorem]{Remark}

\newcommand{\R}{\mathbb{R}}
\newcommand{\Z}{\mathbb{Z}}
\newcommand{\N}{\mathbb{N}}
\newcommand{\C}{\mathbb{C}}
\newcommand{\Q}{\mathbb{Q}}
\newcommand{\F}{\mathbb{F}}
\newcommand{\Sab}{\operatorname{Sab}}
\newcommand{\Aut}{\operatorname{Aut}}
\newcommand{\PSL}{\operatorname{PSL}}
\newcommand{\PGL}{\operatorname{PGL}}
\newcommand{\PSU}{\operatorname{PSU}}
\newcommand{\SU}{\operatorname{SU}}
\newcommand{\GL}{\operatorname{GL}}
\newcommand{\SL}{\operatorname{SL}}
\newcommand{\Ind}{\operatorname{Ind}}
\newcommand{\End}{\operatorname{End}}
\newcommand{\Spec}{\operatorname{Spec}}
\newcommand{\tr}{\operatorname{tr}}
\newcommand{\Frob}{\operatorname{Frob}}
\newcommand{\GH}{\mathrm{GH}}

\title{\textbf{Beyond Me:\\
From Graph Eigenvalues to the Riemann Hypothesis\\
via the Field with One Element}}
\author{Robin Langer and Claude Code}
\date{August 2026}

\begin{document}
\maketitle

\begin{abstract}
We trace a chain of mathematical connections from the concrete to the
conjectural: from the computed eigenvalues of Sabidussi graphs arising
from groups of Lie type, through the Ihara zeta function and the Weil
conjectures, to the Riemann hypothesis and the hypothetical field with
one element~$\F_1$. Along the way we give self-contained definitions
of all objects involved --- Sabidussi graphs, Iwahori--Hecke algebras,
the Ihara zeta function, the Weil zeta function, and the Riemann zeta
function --- and explain why the Ramanujan property of a graph, the
Riemann hypothesis for curves over finite fields, and the classical
Riemann hypothesis are, at bottom, the same statement: that certain
``Frobenius eigenvalues'' have the correct absolute value.

The spectral data computed in the companion papers of this programme
--- all-integer spectra from $G_2(2)$ at $q = 2$, irrational spectra
from $\PSL(2,13)$ at $q = 13$, universal Ramanujan across 30+ graphs
--- provides concrete, verifiable evidence that the $\F_1$-framework
captures something real. Whether this finite evidence points toward
the infinite remains, honestly, beyond us.
\end{abstract}

\tableofcontents

\bigskip

%% =====================================================================
%%
%%  PART I: THE OBJECTS
%%
%% =====================================================================

\part*{Part I: The Objects}
\addcontentsline{toc}{part}{Part I: The Objects}

%% =====================================================================
\section{Sabidussi Graphs}
\label{sec:sabidussi}
%% =====================================================================

\begin{tcolorbox}[breakable, colback=green!8, colframe=green!8,
  boxrule=0pt, sharp corners]
\begin{definition}[Sabidussi graph]
\label{def:sabidussi}
Let $G$ be a finite group, $H \leq G$ a subgroup, and $a \in G$ an
element satisfying $a^{-1} \in HaH$ (equivalently, the double coset
$HaH$ is closed under inversion: $HaH = Ha^{-1}H$). The
\emph{Sabidussi graph} (or \emph{coset graph})
$\Sab(G, H, HaH)$ is the simple graph defined by:
\begin{itemize}
\item \textbf{Vertices}: the set $G/H$ of left cosets of $H$ in $G$.
  Hence $|V| = [G : H]$.
\item \textbf{Edges}: two cosets $xH$ and $yH$ are adjacent if and
  only if $x^{-1}y \in HaH$, i.e.\ the relative position of the two
  cosets lies in the double coset $HaH$.
\end{itemize}
The graph is $d$-regular with $d = |HaH|/|H|$. The group $G$ acts on
$G/H$ by left multiplication, and this action is a group of
automorphisms of $\Sab(G, H, HaH)$.
\end{definition}
\end{tcolorbox}

\begin{tcolorbox}[colback=yellow!12, colframe=yellow!12,
  boxrule=0pt, sharp corners]
\textbf{Arc-transitivity.}
A graph is \emph{arc-transitive} if its automorphism group acts
transitively on ordered pairs of adjacent vertices (arcs). By
Sabidussi's theorem~\cite{Sabidussi1964}, a graph is arc-transitive
if and only if it can be written as $\Sab(G, H, HaH)$ where $HaH$ is
a \emph{single} double coset. When $|HaH|/|H| = 3$, the graph is
cubic (3-regular) and arc-transitive.
\end{tcolorbox}

\begin{tcolorbox}[breakable, colback=black!5, colframe=black!5,
  boxrule=0pt, sharp corners]
\textbf{Examples from the programme.}

\medskip
\textbf{The Langer graph.}
$\Sab(G_2(2), H_{192})$, where $G_2(2)$ is the automorphism group
of the split Cayley hexagon $\GH(2,2)$ (order $12096$) and $H_{192}$
is the stabiliser of a point (order $192$). This gives $63$ vertices,
degree~$6$.

\emph{Spectrum}: $\{6^1,\; 3^{21},\; (-1)^{27},\; (-3)^{14}\}$ ---
all integers.

\medskip
\textbf{The Zhou graph.}
$\Sab(\PSL(2,13), D_{12})$, where $\PSL(2,13)$ has order $1092$ and
$D_{12}$ is a dihedral subgroup of order~$12$. This gives $91$
vertices, degree~$6$.

\emph{Spectrum}: contains $\frac{3 \pm \sqrt{17}}{2}$ (each with
multiplicity~$14$), $\pm 2\sqrt{2}$ (each with multiplicity~$13$),
and three roots of $x^3 + 4x^2 + 3x - 1 = 0$ (each with
multiplicity~$12$). The splitting field is
$\Q(\sqrt{2}, \sqrt{17}, \alpha)$ where $\alpha$ is a cubic
algebraic number.
\end{tcolorbox}

\begin{remark}
The contrast is striking. Both graphs are $6$-regular Sabidussi graphs
of groups of Lie type. The Langer graph, built from $G_2(2)$ over
$\F_2$, has an all-integer spectrum. The Zhou graph, built from
$\PSL(2,13)$ over $\F_{13}$, has eigenvalues involving $\sqrt{17}$
and a cubic irrationality. This difference --- rational at small $q$,
irrational at large $q$ --- is a central theme of Part~II.
\end{remark}


%% =====================================================================
\section{The Iwahori--Hecke Algebra}
\label{sec:hecke}
%% =====================================================================

\subsection{The abstract Hecke algebra of a Coxeter group}

\begin{tcolorbox}[breakable, colback=green!8, colframe=green!8,
  boxrule=0pt, sharp corners]
\begin{definition}[Coxeter group]
\label{def:coxeter}
A \emph{Coxeter group} is a group $W$ with a distinguished set of
generators $S = \{s_1, \ldots, s_n\}$ and a presentation
\[
  W = \langle\, s_1, \ldots, s_n \;\mid\;
  (s_i s_j)^{m_{ij}} = e \;\rangle
\]
where $m_{ii} = 1$ (each generator is an involution) and
$m_{ij} = m_{ji} \in \{2, 3, 4, \ldots, \infty\}$ for $i \neq j$.
The matrix $(m_{ij})$ is the \emph{Coxeter matrix}. When
$m_{ij} = \infty$, no relation is imposed on $s_i s_j$.
\end{definition}
\end{tcolorbox}

\begin{tcolorbox}[breakable, colback=green!8, colframe=green!8,
  boxrule=0pt, sharp corners]
\begin{definition}[Iwahori--Hecke algebra]
\label{def:hecke}
Let $(W, S)$ be a Coxeter group with $S = \{s_1, \ldots, s_n\}$, and
let $q$ be an indeterminate (or a nonzero element of a field). The
\emph{Iwahori--Hecke algebra} $\mathcal{H}_q(W)$ is the associative
$\Q(q)$-algebra with basis $\{T_w : w \in W\}$ (one basis element
per group element) and relations:
\begin{enumerate}
\item \textbf{Quadratic relation}: for each generator $s \in S$,
\begin{equation}
\label{eq:hecke-quad}
  T_s^2 = (q - 1)\, T_s + q\, T_e,
\end{equation}
where $T_e$ is the identity element.

\item \textbf{Braid relations}: if $m_{ij} < \infty$, then
\[
  \underbrace{T_{s_i} T_{s_j} T_{s_i} \cdots}_{m_{ij} \text{ factors}}
  = \underbrace{T_{s_j} T_{s_i} T_{s_j} \cdots}_{m_{ij} \text{ factors}}.
\]

\item \textbf{Multiplicativity on reduced words}: if
$\ell(ww') = \ell(w) + \ell(w')$ (where $\ell$ denotes the length
function on $W$), then $T_{ww'} = T_w T_{w'}$.
\end{enumerate}

As a $\Q(q)$-module, $\mathcal{H}_q(W) \cong \Q(q)^{|W|}$: it has
the same dimension as the group algebra.
\end{definition}
\end{tcolorbox}

\begin{tcolorbox}[breakable, colback=yellow!12, colframe=yellow!12,
  boxrule=0pt, sharp corners]
\textbf{The two key specialisations.}

\begin{itemize}
\item \textbf{At $q = 1$}: the quadratic relation becomes
  $T_s^2 = T_e$, recovering the defining relation $s^2 = e$ of the
  Coxeter group. Thus $\mathcal{H}_1(W) \cong \C[W]$, the group
  algebra of~$W$.

\item \textbf{At $q = p^k$ (a prime power)}: there exists a finite
  group of Lie type $G(\F_q)$ --- for instance
  $\GL_n(\F_q)$, $\SL_n(\F_q)$, or the Chevalley group of type $W$
  over $\F_q$ --- with a Borel subgroup $B$ such that
  \[
    \mathcal{H}_q(W)
    \;\cong\;
    \End_G\!\left(\Ind_B^G \mathbf{1}\right)
    \;=\;
    e_B \cdot \C[G] \cdot e_B
  \]
  where $e_B = \frac{1}{|B|}\sum_{b \in B} b$ is the idempotent
  averaging over~$B$. This is the \emph{Hecke algebra of the pair
  $(G, B)$}: the algebra of $G$-endomorphisms of the representation
  induced from the trivial character of $B$.
\end{itemize}
\end{tcolorbox}

\subsection{The Hecke algebra of a Gelfand pair}

\begin{tcolorbox}[breakable, colback=green!8, colframe=green!8,
  boxrule=0pt, sharp corners]
\begin{definition}[Hecke algebra of a pair $(G, H)$]
\label{def:hecke-pair}
For a finite group $G$ and a subgroup $H$, the \emph{Hecke algebra}
$\mathcal{H}(G, H)$ is the algebra of $H$-bi-invariant functions on
$G$ under convolution:
\[
  \mathcal{H}(G, H) = \End_G\!\left(\Ind_H^G \mathbf{1}\right)
  \;=\; \{f : G \to \C \;\mid\; f(h_1 g h_2) = f(g)
  \text{ for all } h_1, h_2 \in H\}.
\]
As a vector space, $\mathcal{H}(G, H)$ has dimension equal to the
number of $(H, H)$-double cosets in~$G$. The multiplication is
convolution:
\[
  (f_1 * f_2)(g) = \frac{1}{|H|} \sum_{x \in G} f_1(x)\, f_2(x^{-1}g).
\]
\end{definition}
\end{tcolorbox}

\begin{tcolorbox}[breakable, colback=yellow!12, colframe=yellow!12,
  boxrule=0pt, sharp corners]
\textbf{Connection to eigenvalues.}
The eigenvalues of the Sabidussi graph $\Sab(G, H, HaH)$ are the
values of the \emph{zonal spherical functions} of the pair $(G, H)$
evaluated at the double coset $HaH$. A zonal spherical function
$\omega : G \to \C$ is an $H$-bi-invariant function satisfying the
product formula
\[
  \frac{1}{|H|} \sum_{h \in H} \omega(xhy) = \omega(x)\,\omega(y)
  \qquad \text{for all } x, y \in G.
\]
These are precisely the algebra homomorphisms
$\mathcal{H}(G, H) \to \C$, i.e.\ the \emph{characters} of the
Hecke algebra. Each character $\chi$ yields an eigenvalue
$\lambda_\chi = |HaH| \cdot \chi(T_{HaH}) / |H|$ of the adjacency
matrix.

When $(G, H)$ is a \emph{Gelfand pair} (i.e.\ $\mathcal{H}(G, H)$
is commutative), every eigenspace of the adjacency matrix corresponds
to a unique irreducible constituent of $\Ind_H^G \mathbf{1}$, and
the eigenvalue multiplicities equal the dimensions of those
irreducible representations.
\end{tcolorbox}


%% =====================================================================
\section{The Ihara Zeta Function}
\label{sec:ihara}
%% =====================================================================

\begin{tcolorbox}[breakable, colback=green!8, colframe=green!8,
  boxrule=0pt, sharp corners]
\begin{definition}[Prime cycle]
\label{def:prime-cycle}
Let $X$ be a finite connected graph. A \emph{closed walk} of
length~$n$ is a sequence of vertices $v_0, v_1, \ldots, v_n$ with
$v_0 = v_n$ and each $v_i v_{i+1}$ an edge. A closed walk is:
\begin{itemize}
\item \emph{Backtrackless} if $v_{i+1} \neq v_{i-1}$ for all $i$
  (it never immediately reverses).
\item \emph{Tailless} if additionally $v_1 \neq v_{n-1}$ (the
  ``start'' and ``end'' directions differ).
\item \emph{Prime} if it is not a proper power of a shorter closed
  walk: there is no closed walk of length $m < n$ whose $k$-fold
  repetition ($n = km$) gives the original walk.
\end{itemize}
Two prime cycles are \emph{equivalent} if one is a cyclic rotation
of the other: $(v_0, v_1, \ldots, v_{n-1})$ and
$(v_j, v_{j+1}, \ldots, v_{j-1})$ for some~$j$.
\end{definition}
\end{tcolorbox}

\begin{tcolorbox}[breakable, colback=green!8, colframe=green!8,
  boxrule=0pt, sharp corners]
\begin{definition}[Ihara zeta function]
\label{def:ihara}
The \emph{Ihara zeta function} of a finite connected graph $X$ is the
formal power series (in fact, a rational function) defined by the
Euler product over equivalence classes $[C]$ of prime cycles:
\begin{equation}
\label{eq:ihara-euler}
  Z_X(u) = \prod_{[C]} \left(1 - u^{\nu(C)}\right)^{-1},
\end{equation}
where $\nu(C)$ denotes the length of~$C$.
\end{definition}
\end{tcolorbox}

\begin{tcolorbox}[breakable, colback=black!5, colframe=black!5,
  boxrule=0pt, sharp corners]
\begin{theorem}[Ihara determinant formula; Bass \cite{Bass1992},
Hashimoto \cite{Hashimoto1989}]
\label{thm:ihara-det}
Let $X$ be a $(q+1)$-regular graph on $n$ vertices with adjacency
matrix~$A$. Then
\begin{equation}
\label{eq:ihara-det}
  Z_X(u)^{-1}
  = (1 - u^2)^{r-1} \cdot \det\!\left(I - Au + qu^2 I\right),
\end{equation}
where $r = n(q+1)/2 - n + 1 = n(q-1)/2 + 1$ is the rank of the
fundamental group of~$X$ (equivalently, the number of independent
cycles: $r = |E| - |V| + 1$).
\end{theorem}
\end{tcolorbox}

\begin{proof}[Sketch]
Define the \emph{edge matrix} (or Hashimoto matrix) $T$ on directed
edges $\vec{E}$ by $T_{(e_1, e_2)} = 1$ if $e_1$ feeds into $e_2$
without backtracking, and $0$ otherwise. Then
$Z_X(u)^{-1} = \det(I - uT)$. The block structure of $T$ (with
blocks corresponding to edges starting and ending at each vertex)
allows the determinant to be factored into the expression above
involving the adjacency matrix.
\end{proof}

\begin{tcolorbox}[breakable, colback=black!5, colframe=black!5,
  boxrule=0pt, sharp corners]
\begin{proposition}[Functional equation]
\label{prop:ihara-func}
For a $(q+1)$-regular graph on $n$ vertices, the Ihara zeta function
satisfies
\[
  Z_X(1/(qu))
  = (-1)^{n(q-1)/2}\, q^{n(q-1)/2}\, u^{n(q-1)}\, Z_X(u)
  \cdot (\text{elementary factors}).
\]
This is the analogue of the functional equation of classical zeta
functions: it relates $Z_X(u)$ at $u$ to $Z_X$ at the ``dual point''
$1/(qu)$.
\end{proposition}
\end{tcolorbox}

\begin{tcolorbox}[breakable, colback=yellow!12, colframe=yellow!12,
  boxrule=0pt, sharp corners]
\textbf{The Ramanujan criterion.}
The \emph{non-trivial} part of the Ihara zeta function is
$\det(I - Au + qu^2 I)$. Its zeros as a polynomial in $u$ are
determined by the eigenvalues of $A$: if $\lambda$ is an eigenvalue,
the corresponding factor is $1 - \lambda u + qu^2$, with roots
\[
  u = \frac{\lambda \pm \sqrt{\lambda^2 - 4q}}{2q}.
\]
\begin{theorem}[Ramanujan $\Leftrightarrow$ graph RH]
\label{thm:ramanujan-rh}
A $(q+1)$-regular graph $X$ is \emph{Ramanujan} if and only if every
non-trivial zero of $Z_X(u)^{-1}$ satisfies $|u| = q^{-1/2}$.
Equivalently:
\[
  X \text{ is Ramanujan}
  \;\Longleftrightarrow\;
  |\lambda| \leq 2\sqrt{q} \text{ for all non-trivial } \lambda
  \;\Longleftrightarrow\;
  Z_X \text{ satisfies RH}.
\]
\end{theorem}
\end{tcolorbox}


%% =====================================================================
\section{The Weil Conjectures}
\label{sec:weil}
%% =====================================================================

\begin{tcolorbox}[breakable, colback=green!8, colframe=green!8,
  boxrule=0pt, sharp corners]
\begin{definition}[Zeta function of a variety over $\F_q$]
\label{def:weil-zeta}
Let $X$ be a smooth projective algebraic variety of dimension $d$
defined over the finite field $\F_q$. For each $n \geq 1$, let
$N_n = |X(\F_{q^n})|$ denote the number of $\F_{q^n}$-rational
points. The \emph{Hasse--Weil zeta function} is
\begin{equation}
\label{eq:weil-zeta}
  Z(X/\F_q,\, t) = \exp\!\left(\sum_{n=1}^{\infty}
  N_n \frac{t^n}{n}\right).
\end{equation}
\end{definition}
\end{tcolorbox}

\begin{tcolorbox}[breakable, colback=black!5, colframe=black!5,
  boxrule=0pt, sharp corners]
\textbf{The Weil conjectures} (now theorems). Let $X/\F_q$ be a smooth
projective variety of dimension~$d$.

\begin{enumerate}
\item \textbf{Rationality} (Dwork, 1960~\cite{Dwork1960}):
$Z(X/\F_q, t)$ is a rational function of~$t$. Specifically,
\[
  Z(X/\F_q, t)
  = \frac{P_1(t)\, P_3(t) \cdots P_{2d-1}(t)}
         {P_0(t)\, P_2(t) \cdots P_{2d}(t)},
\]
where each $P_i(t) \in \Z[t]$ with $P_0(t) = 1 - t$ and
$P_{2d}(t) = 1 - q^d t$.

\item \textbf{Functional equation} (Grothendieck, via
$\ell$-adic cohomology):
\[
  Z(X/\F_q,\, 1/(q^d t))
  = \pm\, q^{d\chi/2}\, t^{\chi}\, Z(X/\F_q,\, t),
\]
where $\chi = \sum (-1)^i \deg P_i$ is the Euler characteristic.

\item \textbf{The Riemann hypothesis} (Deligne, 1974~\cite{Deligne1974}):
The reciprocal roots of $P_i(t)$ --- i.e.\ the algebraic numbers
$\alpha_{i,j}$ such that
$P_i(t) = \prod_j (1 - \alpha_{i,j} t)$ --- have absolute value
\begin{equation}
\label{eq:deligne-rh}
  |\alpha_{i,j}| = q^{i/2}.
\end{equation}
In particular, the zeros of $Z$ (from odd $P_i$) and poles (from even
$P_i$) lie on circles of radius $q^{-i/2}$ in the $t$-plane.
\end{enumerate}
\end{tcolorbox}

\begin{tcolorbox}[colback=yellow!12, colframe=yellow!12,
  boxrule=0pt, sharp corners]
\textbf{The curve case.} For a smooth projective curve $C/\F_q$ of
genus~$g$, the zeta function simplifies to
\[
  Z(C/\F_q, t) = \frac{P_1(t)}{(1-t)(1-qt)},
\]
where $P_1(t) = \prod_{j=1}^{2g}(1 - \alpha_j t)$ with
$|\alpha_j| = q^{1/2}$. The last statement \emph{is} the Riemann
hypothesis for curves over finite fields, first proved by Weil
(1948~\cite{Weil1948}) for curves and by Deligne
(1974~\cite{Deligne1974}) in general.
\end{tcolorbox}


%% =====================================================================
\section{The Riemann Zeta Function}
\label{sec:riemann}
%% =====================================================================

\begin{tcolorbox}[breakable, colback=green!8, colframe=green!8,
  boxrule=0pt, sharp corners]
\begin{definition}[Riemann zeta function]
\label{def:riemann-zeta}
The \emph{Riemann zeta function} is defined for $\operatorname{Re}(s) > 1$ by
\begin{equation}
\label{eq:riemann-def}
  \zeta(s) = \sum_{n=1}^{\infty} n^{-s}
  = \prod_{p \text{ prime}} \left(1 - p^{-s}\right)^{-1}.
\end{equation}
The Euler product encodes the fundamental theorem of arithmetic:
every positive integer has a unique prime factorisation.
\end{definition}
\end{tcolorbox}

\begin{tcolorbox}[breakable, colback=black!5, colframe=black!5,
  boxrule=0pt, sharp corners]
\textbf{Analytic continuation.}
The function $\zeta(s)$ extends to a meromorphic function on all of
$\C$, with a single simple pole at $s = 1$ with residue~$1$. The
extension can be obtained via the integral representation
\[
  \zeta(s) = \frac{1}{\Gamma(s)} \int_0^{\infty}
  \frac{x^{s-1}}{e^x - 1}\, dx
  \qquad (\operatorname{Re}(s) > 1)
\]
and analytic continuation of the right-hand side.

\medskip
\textbf{Trivial zeros.}
$\zeta(s) = 0$ at $s = -2, -4, -6, \ldots$ (the negative even
integers). These are called \emph{trivial zeros}.
\end{tcolorbox}

\begin{tcolorbox}[breakable, colback=green!8, colframe=green!8,
  boxrule=0pt, sharp corners]
\begin{definition}[Completed zeta function]
\label{def:xi}
The \emph{completed zeta function} is
\begin{equation}
\label{eq:xi}
  \xi(s) = \pi^{-s/2}\, \Gamma(s/2)\, \zeta(s),
\end{equation}
where $\Gamma$ is the gamma function. It satisfies the functional
equation
\begin{equation}
\label{eq:func-eq}
  \xi(s) = \xi(1 - s).
\end{equation}
The factor $\pi^{-s/2}\, \Gamma(s/2)$ is the \emph{Euler factor at
infinity} --- it is the local $L$-factor at the archimedean place of
$\Q$. Just as each prime $p$ contributes a factor
$(1 - p^{-s})^{-1}$ to the Euler product, the ``prime at infinity''
contributes $\pi^{-s/2}\, \Gamma(s/2)$.
\end{definition}
\end{tcolorbox}

\begin{tcolorbox}[colback=blue!8, colframe=blue!8,
  boxrule=0pt, sharp corners]
\textbf{The Riemann Hypothesis (1859).}
Every non-trivial zero of $\zeta(s)$ has real part $\frac{1}{2}$:
\[
  \zeta(s) = 0,\; s \notin \{-2, -4, -6, \ldots\}
  \;\;\Longrightarrow\;\;
  \operatorname{Re}(s) = \tfrac{1}{2}.
\]
Equivalently, the non-trivial zeros of $\xi(s)$ all lie on the
\emph{critical line} $\operatorname{Re}(s) = \frac{1}{2}$.
\end{tcolorbox}


%% =====================================================================
\section{The Ramanujan--Petersson Conjecture}
\label{sec:ramanujan-petersson}
%% =====================================================================

\begin{tcolorbox}[breakable, colback=green!8, colframe=green!8,
  boxrule=0pt, sharp corners]
\begin{definition}[Modular forms and Hecke operators]
\label{def:modular}
A \emph{modular form} of weight $k$ and level $N$ is a holomorphic
function $f : \mathfrak{H} \to \C$ on the upper half-plane
$\mathfrak{H} = \{\tau \in \C : \operatorname{Im}(\tau) > 0\}$
satisfying:
\begin{enumerate}
\item \textbf{Modularity}: for all
$\bigl(\begin{smallmatrix} a & b \\ c & d \end{smallmatrix}\bigr)
\in \Gamma_0(N) \leq \SL_2(\Z)$,
\[
  f\!\left(\frac{a\tau + b}{c\tau + d}\right)
  = (c\tau + d)^k\, f(\tau).
\]
\item \textbf{Holomorphicity at cusps}: $f$ extends holomorphically
  to the cusps of $\Gamma_0(N) \backslash \mathfrak{H}$.
\end{enumerate}
A modular form that vanishes at all cusps is a \emph{cusp form}.
Every cusp form has a Fourier expansion
$f(\tau) = \sum_{n=1}^{\infty} a_n\, e^{2\pi i n \tau}$.

The \emph{Hecke operators} $T_p$ (for primes $p \nmid N$) act on
the space of cusp forms by
\[
  (T_p f)(\tau) = p^{k-1} f(p\tau)
  + \frac{1}{p} \sum_{j=0}^{p-1} f\!\left(\frac{\tau + j}{p}\right).
\]
A \emph{Hecke eigenform} is a cusp form that is a simultaneous
eigenvector of all Hecke operators: $T_p f = a_p f$.
\end{definition}
\end{tcolorbox}

\begin{tcolorbox}[breakable, colback=black!5, colframe=black!5,
  boxrule=0pt, sharp corners]
\begin{theorem}[Ramanujan--Petersson conjecture; proved by
Deligne \cite{Deligne1974}]
\label{thm:ramanujan-petersson}
Let $f = \sum a_n q^n$ be a normalised Hecke eigenform of weight $k$
for $\Gamma_0(N)$. Then for every prime $p \nmid N$,
\begin{equation}
\label{eq:rp-bound}
  |a_p| \leq 2\, p^{(k-1)/2}.
\end{equation}
\end{theorem}
\end{tcolorbox}

\begin{tcolorbox}[breakable, colback=yellow!12, colframe=yellow!12,
  boxrule=0pt, sharp corners]
\textbf{The connection to Ramanujan graphs.}
A $(q+1)$-regular graph is Ramanujan when its non-trivial eigenvalues
satisfy $|\lambda| \leq 2\sqrt{q}$. This is \emph{exactly} the
Ramanujan--Petersson bound~\eqref{eq:rp-bound} at weight $k = 2$:
\[
  2\,p^{(k-1)/2}\big|_{k=2} = 2\sqrt{p} = 2\sqrt{q}.
\]
The name ``Ramanujan graph'' was chosen by Lubotzky, Phillips, and
Sarnak~\cite{LPS} precisely because of this connection. Their
$(p+1)$-regular Cayley graphs of $\PSL(2, \Z/q\Z)$ are Ramanujan as
a direct consequence of Deligne's theorem.

Deligne proved the Ramanujan--Petersson conjecture by showing that the
eigenvalues $a_p$ are traces of the Frobenius endomorphism on the
$\ell$-adic cohomology $H^1$ of a certain algebraic variety (the
universal elliptic curve over the modular curve), and then applying
his proof of the Weil conjecture~\eqref{eq:deligne-rh} at $i = 1$.
\end{tcolorbox}


%% =====================================================================
\section{The Field with One Element}
\label{sec:f1}
%% =====================================================================

\subsection{Tits's observation}

\begin{tcolorbox}[breakable, colback=yellow!12, colframe=yellow!12,
  boxrule=0pt, sharp corners]
In 1957, Jacques Tits~\cite{Tits1957} observed a remarkable pattern.
Many formulas in the theory of algebraic groups over $\F_q$ have
well-defined limits as $q \to 1$, and these limits recover classical
combinatorial objects:

\medskip
\begin{center}
\renewcommand{\arraystretch}{1.4}
\begin{tabular}{lcc}
\toprule
Object at $q$ & Formula & Limit as $q \to 1$ \\
\midrule
$q$-integer $[n]_q$ & $\dfrac{q^n - 1}{q - 1}$ & $n$ \\[6pt]
$q$-factorial $[n]_q!$ & $\displaystyle\prod_{k=1}^n [k]_q$ & $n!$ \\[6pt]
Gaussian binomial $\displaystyle\binom{n}{k}_q$
  & $\dfrac{[n]_q!}{[k]_q!\,[n-k]_q!}$ & $\displaystyle\binom{n}{k}$ \\[6pt]
$|\GL_n(\F_q)|$ & $\displaystyle\prod_{k=0}^{n-1}(q^n - q^k)$ & $n!$ \\[6pt]
$|\mathbb{P}^{n-1}(\F_q)|$ & $[n]_q = \dfrac{q^n-1}{q-1}$ & $n$ \\
\bottomrule
\end{tabular}
\end{center}

\medskip
This suggests the existence of a ``field with one element''~$\F_1$
such that:
\begin{itemize}
\item $\GL_n(\F_1) = S_n$ (the symmetric group).
\item $\mathbb{P}^{n-1}(\F_1) = \{1, 2, \ldots, n\}$
  (a bare finite set).
\item The Weyl group $W$ of a reductive group $G$ is ``$G(\F_1)$''.
\end{itemize}

This is exactly what the Hecke algebra encodes: $\mathcal{H}_q(W)$ at
$q = 1$ gives $\C[W]$, and at $q = p^k$ gives the Hecke algebra of
the Lie-type group over $\F_q$. The Hecke algebra is a
\emph{deformation from $\F_1$ to $\F_q$}.
\end{tcolorbox}

\subsection{The dream}

\begin{tcolorbox}[breakable, colback=blue!8, colframe=blue!8,
  boxrule=0pt, sharp corners]
\textbf{The $\F_1$ programme.}
The most ambitious goal of $\F_1$-geometry is to prove the Riemann
hypothesis by analogy with Deligne's proof of the Weil conjectures.
The strategy:

\begin{enumerate}
\item View $\Spec(\Z)$ as a ``curve over $\F_1$''. The closed points
  of $\Spec(\Z)$ are the primes $p$; each has a ``residue field''
  $\F_p$ of ``size'' $p$. The ``generic point'' is $\Spec(\Q)$.
  The archimedean absolute value $|\cdot|_\infty$ plays the role of
  a ``point at infinity''.

\item Construct a Weil-type cohomology theory for varieties over
  $\F_1$. For curves over $\F_q$, Deligne used $\ell$-adic
  \'etale cohomology $H^i(X \times_{\F_q} \overline{\F_q}, \Q_\ell)$,
  on which the geometric Frobenius acts with eigenvalues of absolute
  value $q^{i/2}$. The ``Frobenius at $p$'' on $\Spec(\Z)$ is the
  map $x \mapsto x^p$ on the ``fibre'' $\Spec(\F_p)$.

\item Show that the zeros of $\zeta(s)$ are eigenvalues of a
  Frobenius-like operator on this cohomology, and that the ``Weil
  conjectures for $\F_1$'' force them to lie on $\operatorname{Re}(s) = 1/2$.
\end{enumerate}

\medskip
Several approaches have been proposed:
\begin{itemize}
\item \textbf{Soul\'e} (1999~\cite{Soule2004}): defined $\F_1$-schemes
  via a functor category; the ``zeta function'' of $\Spec(\Z)$ in his
  framework involves the completed Riemann zeta function $\xi(s)$.
\item \textbf{Connes--Consani} (2010--\cite{ConnesConsani2011}):
  introduced the \emph{arithmetic site}, a topos built from the
  multiplicative monoid $\N^\times$, and showed that its
  ``characteristic 1'' geometry recovers the explicit formulas of
  prime number theory.
\item \textbf{Borger} (2009~\cite{Borger2009}): proposed that
  ``$\F_1$-geometry \emph{is} $\Lambda$-ring geometry.'' A
  $\Lambda$-ring is a ring equipped with commuting Frobenius lifts
  $\psi_p$ for all primes $p$; $\Z$ with $\psi_p(n) = n$ is the
  ``final object'' over $\F_1$.
\item \textbf{Lorscheid} (2012~\cite{Lorscheid2018}): introduced
  \emph{blueprints} (monoids with a partial addition) as a foundation
  for $\F_1$-geometry that recovers both Soul\'e's and Connes's
  frameworks.
\end{itemize}
\end{tcolorbox}

\begin{tcolorbox}[colback=yellow!12, colframe=yellow!12,
  boxrule=0pt, sharp corners]
\textbf{The key idea.}
If $\Spec(\Z)$ can be treated as a curve over $\F_1$, and if a
cohomology theory can be built for which the ``Frobenius eigenvalues''
are the zeros of $\zeta(s)$, then the classical Riemann hypothesis
would follow from exactly the same formal argument that Deligne used
for curves over $\F_q$: the eigenvalues of Frobenius on $H^1$ have
absolute value $q^{1/2}$, which at $q = 1$ becomes
$|$eigenvalue$| = 1$, which corresponds to
$\operatorname{Re}(s) = 1/2$.
\end{tcolorbox}


%% =====================================================================
%%
%%  PART II: THE CONNECTIONS
%%
%% =====================================================================

\part*{Part II: The Connections}
\addcontentsline{toc}{part}{Part II: The Connections}

%% =====================================================================
\section{Why Our Graphs are Ramanujan}
\label{sec:why-ramanujan}
%% =====================================================================

\begin{tcolorbox}[breakable, colback=black!5, colframe=black!5,
  boxrule=0pt, sharp corners]
\textbf{The argument, made precise.}
Let $G = G(\F_q)$ be a finite group of Lie type over $\F_q$, and let
$X = \Sab(G, H, HaH)$ be a Sabidussi graph. The adjacency matrix $A$
of $X$ acts on the permutation module $\C[G/H] = \Ind_H^G \mathbf{1}$.
Decompose this into irreducibles:
\[
  \Ind_H^G \mathbf{1} = \bigoplus_{\pi \in \widehat{G}^H}
  m_\pi \cdot \pi,
\]
where $\widehat{G}^H = \{\pi \in \widehat{G} : \pi^H \neq 0\}$ is
the set of irreducible representations having a nonzero $H$-fixed
vector, and $m_\pi = \dim(\pi^H)$. On each copy of $\pi$, the
adjacency operator acts by the scalar
\[
  \lambda_\pi = \frac{|HaH|}{|H|} \cdot \omega_\pi(a)
  = d \cdot \omega_\pi(a),
\]
where $\omega_\pi$ is the zonal spherical function associated to
$\pi$, and $d$ is the degree of the graph.

The trivial representation $\pi = \mathbf{1}$ contributes the
eigenvalue $\lambda_{\mathbf{1}} = d$ (the degree). The remaining
eigenvalues are the \emph{non-trivial} ones. The Ramanujan property
asserts that $|\lambda_\pi| \leq 2\sqrt{d-1}$ for all non-trivial
$\pi$.
\end{tcolorbox}

\begin{tcolorbox}[breakable, colback=yellow!12, colframe=yellow!12,
  boxrule=0pt, sharp corners]
\textbf{The PSL(2,$q$) case.}
For $G = \PSL(2, q)$, the spherical functions are values of matrix
coefficients of the principal series representations. By the
Jacquet--Langlands correspondence (or directly via the Eichler--Shimura
relation), these matrix coefficients are related to Hecke eigenvalues
on modular forms.

Concretely, let $X_0(q)$ be the modular curve of level $q$, viewed as
a curve over $\F_p$ for some auxiliary prime $p \neq q$. The Frobenius
$\Frob_p$ acts on the \'etale cohomology $H^1(X_0(q) \times
\overline{\F_p}, \Q_\ell)$, and its eigenvalues $\alpha_p$ satisfy
$|\alpha_p| = p^{1/2}$ by Deligne's theorem~\eqref{eq:deligne-rh}.
The Hecke eigenvalue $a_p$ of the associated modular form satisfies
$a_p = \alpha_p + \overline{\alpha}_p$, giving
$|a_p| \leq 2p^{1/2}$ --- the Ramanujan--Petersson bound.

The eigenvalues of $\Sab(\PSL(2,q), H)$ are precisely (up to
normalisation) these Hecke eigenvalues, evaluated at representations
determined by the pair $(G, H)$. The Ramanujan bound for the graph
is the Ramanujan--Petersson bound for the corresponding modular form,
which is Deligne's theorem.
\end{tcolorbox}

\begin{tcolorbox}[breakable, colback=black!5, colframe=black!5,
  boxrule=0pt, sharp corners]
\textbf{The $G_2$ case.}
For $G = G_2(2)$, the connection to automorphic forms is through the
Kim--Shahidi functoriality result~\cite{KimShahidi2002}. The
symmetric-square and exterior-square $L$-functions associated to
$G_2$ automorphic representations can be related to $\GL(7)$
$L$-functions, to which Deligne's theorem applies. However, for the
smallest case $q = 2$, the Ramanujan bound is satisfied for more
elementary reasons: the group $G_2(2)$ has order $12096$, and its
character table (computable by hand or by GAP) directly gives the
eigenvalues as integers $\{6, 3, -1, -3\}$, all satisfying
$|\lambda| \leq 2\sqrt{5} \approx 4.47$.
\end{tcolorbox}


%% =====================================================================
\section{The $q \to 1$ Limit and Spectral Integrality}
\label{sec:q-to-1}
%% =====================================================================

\begin{tcolorbox}[breakable, colback=yellow!12, colframe=yellow!12,
  boxrule=0pt, sharp corners]
\textbf{The Hecke algebra as a bridge.}
Recall the two specialisations of $\mathcal{H}_q(W)$:
\[
  \begin{array}{ccccc}
  \C[W] & \xleftarrow{\;q \to 1\;} & \mathcal{H}_q(W)
  & \xrightarrow{\;q = p^k\;} & \mathcal{H}(G(\F_q), B) \\
  \text{group algebra} & & \text{deformation} & & \text{Lie-type Hecke}
  \end{array}
\]
At $q = 1$, the eigenvalues of $\mathcal{H}_1(W) = \C[W]$ are
character values of the Weyl group $W$. Character values of finite
groups are algebraic integers --- in fact, they lie in the ring of
integers of cyclotomic fields. For the Weyl groups ($S_n$, $B_n$,
$D_n$, etc.), all character values are \emph{rational} integers.

At $q = p^k$, the eigenvalues are values of spherical functions of
$G(\F_q)$. These lie in the algebraic number field generated by the
character values of $G(\F_q)$ and the parameter $q$. As $q$ grows, the
degree of this number field over $\Q$ generically increases.
\end{tcolorbox}

\begin{tcolorbox}[breakable, colback=black!5, colframe=black!5,
  boxrule=0pt, sharp corners]
\textbf{The evidence from our programme.}

\medskip
\begin{center}
\renewcommand{\arraystretch}{1.3}
\begin{tabular}{lcccl}
\toprule
Graph & $q$ & $|V|$ & Spectrum & Splitting field \\
\midrule
$\Sab(\PSL(2,5), S_3)$ & 5 & 10 & $\{3, 1, -2\}$
  & $\Q$ \\
$\Sab(\PSL(2,5), C_3)$ & 5 & 20
  & $\{3, \sqrt{5}, 1, -\sqrt{5}, -2\}$ & $\Q(\sqrt{5})$ \\
$\Sab(G_2(2), H_{192})$ & 2 & 63
  & $\{6, 3, -1, -3\}$ & $\Q$ \\
$\Sab(\PSL(2,7), C_3)$ & 7 & 56
  & $\ni\, 1+\sqrt{2}$ & $\Q(\sqrt{2})$ \\
$\Sab(\PSL(2,13), D_{12})$ & 13 & 91
  & $\ni\, \frac{3+\sqrt{17}}{2}$
  & $\Q(\sqrt{2},\sqrt{17},\alpha)$ \\
$\Sab(\PSL(2,17), S_3)$ & 17 & 408
  & $\ni\, 2.709\ldots$ & Higher degree \\
\bottomrule
\end{tabular}
\end{center}

\medskip
The pattern is unmistakable: at $q = 2$ and $q = 5$ (with the right
stabiliser), the spectrum is rational --- even integral. At $q = 7$
and $q = 13$, square roots appear. At $q = 17$, the algebraic
complexity grows further.
\end{tcolorbox}

\begin{tcolorbox}[colback=yellow!12, colframe=yellow!12,
  boxrule=0pt, sharp corners]
\begin{observation}[Spectral integrality as proximity to $\F_1$]
\label{obs:integrality}
The Langer graph $\Sab(G_2(2), H_{192})$, built from $G_2$ over
$\F_2$, has an all-integer spectrum $\{6^1, 3^{21}, (-1)^{27},
(-3)^{14}\}$. This is because $q = 2$ is ``close'' to $q = 1$ in the
following precise sense: the Hecke algebra $\mathcal{H}_2(W(G_2))$
has structure constants that, when specialised at $q = 2$, produce
rational eigenvalues. The small distance from the $q = 1$ limit
(where eigenvalues are character values of the Weyl group, hence
integers) means the deformation has not yet introduced
irrationality.

The Zhou graph, built from $\PSL(2,13)$ over $\F_{13}$, is ``far''
from $q = 1$. The Hecke algebra at $q = 13$ has genuinely irrational
spherical function values, and the spectrum contains $\sqrt{17}$.

\emph{Spectral integrality is a measure of proximity to $\F_1$.}
\end{observation}
\end{tcolorbox}


%% =====================================================================
\section{The Three Riemann Hypotheses}
\label{sec:three-rh}
%% =====================================================================

\begin{tcolorbox}[breakable, colback=black!5, colframe=black!5,
  boxrule=0pt, sharp corners]
\begin{center}
\renewcommand{\arraystretch}{1.6}
\begin{tabular}{@{}p{2.5cm} p{3.5cm} p{4cm} p{2.5cm}@{}}
\toprule
Setting & Zeta function & Riemann Hypothesis & Status \\
\midrule
Variety $X/\F_q$
  & $Z(X/\F_q, t) = \exp(\sum N_n t^n/n)$
  & Reciprocal roots of $P_i$ have $|\alpha| = q^{i/2}$
  & \textbf{Proved} (Deligne, 1974) \\
$(q{+}1)$-regular graph $X$
  & $Z_X(u) = \prod_{[C]}(1-u^{\nu(C)})^{-1}$
  & Non-trivial poles at $|u| = q^{-1/2}$
  & $\Leftrightarrow$ Ramanujan \\
$\Spec(\Z)$
  & $\zeta(s) = \prod_p (1-p^{-s})^{-1}$
  & $\operatorname{Re}(s) = 1/2$
  & \textbf{Open} \\
``$\Spec(\Z)$ over $\F_1$''
  & ???
  & Would imply classical RH
  & \textbf{Dream} \\
\bottomrule
\end{tabular}
\end{center}
\end{tcolorbox}

\begin{tcolorbox}[breakable, colback=yellow!12, colframe=yellow!12,
  boxrule=0pt, sharp corners]
\textbf{The punchline: all three are the same statement.}

Each row of the table above is an instance of the same assertion:
\emph{the eigenvalues of a Frobenius-like operator have the correct
absolute value.}

\begin{itemize}
\item Over $\F_q$: the Frobenius endomorphism $\Frob_q : x \mapsto
  x^q$ acts on \'etale cohomology $H^i(X \times \overline{\F_q},
  \Q_\ell)$. Its eigenvalues have absolute value $q^{i/2}$.

\item On a graph: the Hashimoto edge operator $T$ (the
  ``non-backtracking random walk'') plays the role of Frobenius.
  Its eigenvalues (equivalently, the adjacency eigenvalues via the
  Ihara formula) satisfy the analogue of RH precisely when the
  graph is Ramanujan.

\item Over $\Z$: the Frobenius at each prime $p$ acts on the
  (hypothetical) cohomology of $\Spec(\Z)$. The zeros of $\zeta(s)$
  are the eigenvalues. The Riemann hypothesis is the statement that
  these eigenvalues have the ``correct'' absolute value.

\item Over $\F_1$: if the $\F_1$-framework can be made rigorous
  enough to support a Weil cohomology theory, then the classical RH
  would follow from the same formal argument Deligne used.
\end{itemize}
\end{tcolorbox}


%% =====================================================================
\section{The Horror of the Real Interval $(0,1)$}
\label{sec:horror}
%% =====================================================================

\begin{tcolorbox}[breakable, colback=yellow!12, colframe=yellow!12,
  boxrule=0pt, sharp corners]
\textbf{The clean world and the wild world.}

The proofs of the Weil conjectures live in the $p$-adic world: the
fields $\Q_p$ with their ultrametric absolute values, totally
disconnected topologies, and the Frobenius endomorphism $x \mapsto
x^p$ acting cleanly on \'etale cohomology. The $p$-adic world is
\emph{discrete}: every triangle is isoceles, open balls are closed,
and Galois theory controls everything.

The Riemann hypothesis lives in the archimedean world: $\R$ and $\C$
with their ordinary absolute values, the topology of the continuum,
and all its attendant pathologies --- non-measurable sets,
Banach--Tarski, the continuum hypothesis. The archimedean world is
\emph{continuous}: the analytic structure of $\zeta(s)$ is governed
by complex analysis, not algebra.
\end{tcolorbox}

\begin{tcolorbox}[breakable, colback=black!5, colframe=black!5,
  boxrule=0pt, sharp corners]
\textbf{The Gamma factor.}
The completed zeta function $\xi(s) = \pi^{-s/2}\Gamma(s/2)\zeta(s)$
includes the factor $\pi^{-s/2}\Gamma(s/2)$, which is the
\emph{Euler factor at infinity}. In the Euler product decomposition
\[
  \xi(s) = \prod_{v} L_v(s),
\]
each finite place $v = p$ contributes $L_p(s) = (1 - p^{-s})^{-1}$,
and the archimedean place $v = \infty$ contributes
$L_\infty(s) = \pi^{-s/2}\Gamma(s/2)$.

The Gamma factor encodes the contribution of the continuous spectrum.
It is where the archimedean place --- the ``prime at infinity'' ---
makes its presence felt. The Weil conjectures handle all the finite
primes. The Riemann hypothesis is, in this sense, a statement about
the single remaining prime: the archimedean one.
\end{tcolorbox}

\begin{tcolorbox}[breakable, colback=yellow!12, colframe=yellow!12,
  boxrule=0pt, sharp corners]
\textbf{$\F_1$ as the threshold.}
The field $\F_1$ sits at $q = 1$: the threshold between the
ultrametric world ($q = p^k$, discrete, algebraic) and the
archimedean world ($q \to \infty$, continuous, analytic). The Hecke
algebra $\mathcal{H}_q(W)$ deforms the group algebra $\C[W]$ (at
$q = 1$, the ``$\F_1$ point'') into the Lie-type Hecke algebra (at
$q = p^k$, the ``$\F_q$ points'').

\medskip
\textbf{The Ramanujan bound as a spectral edge.}
For a $(q+1)$-regular graph, the Ramanujan bound
$|\lambda| \leq 2\sqrt{q}$ is the edge of the
\emph{Plancherel measure} on the unitary dual of the free group
$F_{q+1}$. The Kesten--McKay distribution on $[-2\sqrt{q}, 2\sqrt{q}]$
is the spectral measure of the infinite $(q+1)$-regular tree. A graph
is Ramanujan precisely when its non-trivial spectrum stays within the
support of this measure.

In representation-theoretic terms: the non-trivial eigenvalues come
from \emph{tempered} representations --- those appearing in the
regular representation, whose matrix coefficients are
$L^{2+\epsilon}$. Non-Ramanujan eigenvalues come from
\emph{non-tempered} (complementary series) representations, where
the matrix coefficients grow too fast.

Being Ramanujan means the spectrum has not crossed from the algebraic
world into the analytic wilderness.
\end{tcolorbox}


%% =====================================================================
%%
%%  PART III: SPECULATION
%%
%% =====================================================================

\part*{Part III: Speculation}
\addcontentsline{toc}{part}{Part III: Speculation}


%% =====================================================================
\section{What the Eigenvalue Data Suggests}
\label{sec:data}
%% =====================================================================

\begin{tcolorbox}[breakable, colback=black!5, colframe=black!5,
  boxrule=0pt, sharp corners]
\textbf{Summary of computed spectra.}
Across the companion papers of this programme, we have computed
the complete spectrum of 30+ Sabidussi graphs from groups of Lie type.
The key findings:

\begin{enumerate}
\item \textbf{Universal Ramanujan property.}
Every graph from a simple Lie-type group is Ramanujan for its degree.
This holds across:
\begin{itemize}
\item All 13 Foster census graphs (cubic,
  $2\sqrt{2} \approx 2.828$).
\item All 7 $\PSL(2,q)$ graphs ($q = 5, 7, 11, 13, 17$).
\item All 3 $G_2(2)$ graphs ($6$-regular, $7$-regular, cubic).
\item All 6 cubic residue voltage lifts.
\end{itemize}

\item \textbf{Integrality at small $q$, irrationality at large $q$.}
The splitting field of the characteristic polynomial grows with $q$:
\begin{center}
\small
\begin{tabular}{llll}
\toprule
$q$ & Graph & Splitting field & Irrational eigenvalues \\
\midrule
$2$ & Langer ($G_2$) & $\Q$ & None \\
$2$ & Tutte 12-cage ($G_2$) & $\Q(\sqrt{6})$ & $\sqrt{6}$ \\
$5$ & Petersen ($A_1$) & $\Q$ & None \\
$5$ & Dodecahedron ($A_1$) & $\Q(\sqrt{5})$ & $\sqrt{5}$ \\
$7$ & Klein ($A_1$) & $\Q(\sqrt{2})$ & $1 + \sqrt{2}$ \\
$13$ & Zhou ($A_1$) & $\Q(\sqrt{2},\sqrt{17},\alpha)$
  & $\frac{3+\sqrt{17}}{2}$, $2\sqrt{2}$, $\alpha$ \\
$17$ & Biggs-Smith ($A_1$) & $\Q(\sqrt{2},\ldots)$ & $2.562\ldots$ \\
\bottomrule
\end{tabular}
\end{center}

\item \textbf{Eigenvalue multiplicities from representation theory.}
The multiplicity of each eigenvalue equals the dimension of the
corresponding irreducible constituent of $\Ind_H^G \mathbf{1}$.
For the Langer graph: multiplicities $1, 21, 27, 14$ correspond to
dimensions $1, 21, 27, 14$ of the four irreducible
$G_2(2)$-representations appearing in the permutation character.
\end{enumerate}
\end{tcolorbox}

\begin{tcolorbox}[colback=blue!8, colframe=blue!8,
  boxrule=0pt, sharp corners]
\begin{conjecture}[Spectral complexity growth]
\label{conj:complexity}
The spectral complexity of $\Sab(G, H)$ over $\F_q$, measured by the
degree of the splitting field of the characteristic polynomial of $A$
over $\Q$, is monotonically increasing in $q$ for fixed Lie type and
fixed stabiliser type.

More precisely: let $G_q = G(\F_q)$ be a family of Chevalley groups
of fixed type (e.g.\ $\PSL(2,q)$, or $G_2(q)$), with stabiliser
$H_q$ of fixed type (e.g.\ $C_3$, or $D_{2(q-1)}$, or the point
stabiliser). Let $f_q(x) = \det(xI - A_q)$ be the characteristic
polynomial of the adjacency matrix. Then
\[
  [\Q(\text{roots of } f_q) : \Q]
  \;\leq\;
  [\Q(\text{roots of } f_{q'}) : \Q]
  \qquad \text{whenever } q \leq q'.
\]
\end{conjecture}
\end{tcolorbox}


%% =====================================================================
\section{Toward a Proof?}
\label{sec:toward}
%% =====================================================================

\begin{tcolorbox}[breakable, colback=black!5, colframe=black!5,
  boxrule=0pt, sharp corners]
\textbf{What would be needed.}
To prove the Riemann hypothesis via the $\F_1$ route, one would need:

\begin{enumerate}
\item \textbf{A rigorous definition of $\F_1$-geometry.}
Multiple candidates exist (Soul\'e, Connes--Consani, Borger,
Lorscheid), but none has yet produced the full machinery needed.

\item \textbf{A Weil cohomology for $\Spec(\Z)$ over $\F_1$.}
This would be a cohomology theory $H^i(\Spec(\Z), ?)$ satisfying:
\begin{itemize}
\item $H^0 = \Q$ (one connected component).
\item $H^2 \cong \Q(1)$ (Poincar\'e duality for a ``curve'').
\item $H^1$ has dimension equal to the number of zeros of $\xi(s)$
  on the critical line.
\item A ``Frobenius'' operator $F$ on $H^1$ whose eigenvalues are
  the non-trivial zeros of $\zeta$.
\end{itemize}

\item \textbf{A positivity result.}
The analogue of Deligne's use of the Lefschetz trace formula and the
Hodge index theorem: a positivity argument that forces the
eigenvalues to have the correct absolute value.
\end{enumerate}
\end{tcolorbox}

\begin{tcolorbox}[breakable, colback=black!5, colframe=black!5,
  boxrule=0pt, sharp corners]
\textbf{Connes--Consani: the arithmetic site.}
The most developed approach is that of Connes and Consani. Their
\emph{arithmetic site} $(\widehat{\N^\times}, \Z_{\max})$ is a topos
equipped with a structure sheaf valued in the tropical semifield
$\Z_{\max} = (\Z \cup \{-\infty\}, \max, +)$. The ``points'' of this
site over the ``algebraic closure of $\F_1$'' (which they identify
with $\R_+^\times$) are the ad\`ele classes
$\Q^\times \backslash \mathbb{A}_\Q / \widehat{\Z}^\times$, and the
``scaling action'' of $\R_+^\times$ on these ad\`ele classes plays the
role of the Frobenius.

In this framework, the Weil explicit formulas of prime number theory
(which relate sums over primes to sums over zeros of $\zeta$) become
a Lefschetz trace formula for the ``Frobenius correspondences'' on
the arithmetic site.

\medskip
\textbf{Borger: $\Lambda$-rings.}
Borger's approach identifies $\F_1$-algebras with $\Lambda$-rings ---
commutative rings equipped with commuting Frobenius lifts $\psi_p$
for all primes $p$ (equivalently, a $\Lambda$-ring structure).
The ring $\Z$ with $\psi_p = \operatorname{id}$ is the initial
$\Lambda$-ring, hence ``$\Spec(\Z) = \Spec(\F_1)$.'' In this setup,
the Witt vectors $W(\F_p) = \Z_p$ play the role of the ``residue
field extensions'' of $\F_1$ at the prime $p$.
\end{tcolorbox}

\begin{tcolorbox}[breakable, colback=yellow!12, colframe=yellow!12,
  boxrule=0pt, sharp corners]
\textbf{Why it's ``Beyond Me.''}
The gap between what we can prove and what we dream of proving is
vast. We can compute eigenvalues. We can verify the Ramanujan property.
We can trace the chain of analogies from graph spectra through the
Ihara zeta function and the Weil conjectures to the Riemann
hypothesis. We can observe that spectral integrality at small $q$
reflects proximity to $\F_1$, and that the Hecke algebra
interpolates between the $\F_1$-world and the $\F_q$-world.

But the final step --- constructing the $\F_1$-cohomology of
$\Spec(\Z)$ and proving the positivity result that would force the
zeros of $\zeta$ onto the critical line --- remains beyond the reach
of current mathematics. The machinery that Grothendieck built for the
Weil conjectures (schemes, \'etale cohomology, the six functors) took
decades and hundreds of pages of EGA and SGA. The analogous machinery
for $\F_1$ is still in its infancy.

What we can say is this: the spectral data from our graphs is
\emph{consistent} with the $\F_1$ picture. The Hecke algebra really
does interpolate between the symmetric group (at $q = 1$) and the
Lie-type group (at $q = p^k$). The eigenvalues really are rational at
small $q$ and irrational at large $q$. The Ramanujan property really
does hold universally for Lie-type Sabidussi graphs. And the formal
structure of the Ihara zeta function really does mirror the Hasse--Weil
zeta function, with the Ramanujan condition playing the role of the
Riemann hypothesis.

The graphs are computable. The eigenvalues are verifiable. They all
satisfy RH. The question is whether this finite evidence --- computed
from groups of order $12096$ and fields of size $13$ --- points toward
the infinite.
\end{tcolorbox}

\begin{tcolorbox}[breakable, colback=black!5, colframe=black!5,
  boxrule=0pt, sharp corners]
\textbf{Final observation.}
There is a sense in which the chain of analogies we have traced is
itself the content. The Ramanujan property of a graph, the Weil
conjectures for a variety, and the Riemann hypothesis for the integers
are not merely \emph{analogous} statements --- they are shadows of a
single statement, projected onto different mathematical worlds:

\medskip
\begin{center}
\emph{The eigenvalues of Frobenius have the correct absolute value.}
\end{center}

\medskip
At $q = p^k$, we know this is true (Deligne). At $q = 1$, we
believe it is true (RH). The Hecke algebra, which deforms from $q = 1$
to $q = p^k$, is the bridge. Our graphs are points on this bridge,
close enough to $q = 1$ to compute, far enough from $q = 1$ to
exhibit the phenomena --- irrationality, algebraic number fields,
the growth of spectral complexity --- that make the connection
non-trivial.

Whether the bridge can be crossed remains, for now, beyond us.
\end{tcolorbox}


%% =====================================================================
%%  BIBLIOGRAPHY
%% =====================================================================

\begin{thebibliography}{99}

\bibitem{Bass1992}
H.~Bass,
\emph{The Ihara--Selberg zeta function of a tree lattice},
Internat.\ J.\ Math.\ \textbf{3} (1992), no.~6, 717--797.

\bibitem{Borger2009}
J.~Borger,
\emph{$\Lambda$-rings and the field with one element},
preprint, arXiv:0906.3146, 2009.

\bibitem{ConnesConsani2011}
A.~Connes and C.~Consani,
\emph{On the notion of geometry over $\F_1$},
J.\ Algebraic Geom.\ \textbf{20} (2011), no.~3, 525--557.

\bibitem{Deligne1974}
P.~Deligne,
\emph{La conjecture de Weil.\ I},
Inst.\ Hautes \'Etudes Sci.\ Publ.\ Math.\ \textbf{43} (1974), 273--307.

\bibitem{Dwork1960}
B.~Dwork,
\emph{On the rationality of the zeta function of an algebraic variety},
Amer.\ J.\ Math.\ \textbf{82} (1960), 631--648.

\bibitem{Hashimoto1989}
K.~Hashimoto,
\emph{Zeta functions of finite graphs and representations of $p$-adic groups},
in: Automorphic Forms and Geometry of Arithmetic Varieties,
Adv.\ Stud.\ Pure Math.\ \textbf{15}, Academic Press, 1989, pp.~211--280.

\bibitem{Ihara1966}
Y.~Ihara,
\emph{On discrete subgroups of the two by two projective linear group over
$\mathfrak{p}$-adic fields},
J.\ Math.\ Soc.\ Japan \textbf{18} (1966), 219--235.

\bibitem{KimShahidi2002}
H.~Kim and F.~Shahidi,
\emph{Functorial products for $\GL_2 \times \GL_3$ and the symmetric cube
for $\GL_2$},
Ann.\ of Math.\ (2) \textbf{155} (2002), no.~3, 837--893.

\bibitem{Langer2013}
R.~Langer,
\emph{Symmetric quotients of arc-transitive graphs},
Honours thesis, University of Melbourne, 2013.
arXiv:1306.4798.

\bibitem{Lorscheid2018}
O.~Lorscheid,
\emph{$\F_1$ for everyone},
Jahresber.\ Dtsch.\ Math.-Ver.\ \textbf{120} (2018), 83--116.

\bibitem{LPS}
A.~Lubotzky, R.~Phillips, and P.~Sarnak,
\emph{Ramanujan graphs},
Combinatorica \textbf{8} (1988), no.~3, 261--277.

\bibitem{Sabidussi1964}
G.~Sabidussi,
\emph{Vertex-transitive graphs},
Monatsh.\ Math.\ \textbf{68} (1964), 426--438.

\bibitem{Soule2004}
C.~Soul\'e,
\emph{Les vari\'et\'es sur le corps \`a un \'el\'ement},
Mosc.\ Math.\ J.\ \textbf{4} (2004), no.~1, 217--244.

\bibitem{Tits1957}
J.~Tits,
\emph{Sur les analogues alg\'ebriques des groupes semi-simples complexes},
in: Colloque d'alg\`ebre sup\'erieure, Bruxelles, 1957, pp.~261--289.

\bibitem{Weil1948}
A.~Weil,
\emph{Vari\'et\'es ab\'eliennes et courbes alg\'ebriques},
Actualit\'es Sci.\ Ind., no.~1064, Hermann, Paris, 1948.

\end{thebibliography}

\end{document}
